\documentclass[11pt]{article}
\usepackage[margin=1in]{geometry}
\usepackage[hidelinks]{hyperref}
\usepackage{tabularx}
\usepackage{booktabs}
\usepackage{url}

\title{Specular-Aware Monocular Depth and Clinical Adaptation Guardrails\\
\large Preliminary Findings for Colonoscopy Polyp-Size Research}
\author{Smart Wattanapornmongkol\thanks{\raggedright Handoff contact: student ID 6738249221; phone 0639390653; \href{https://line.me/ti/p/nqlvyY38hD}{LINE}.\\Project repository: \href{https://github.com/smartwhatt/Colon-AF-SfMLearner}{github.com/smartwhatt/Colon-AF-SfMLearner}.}}
\date{23 July 2026}

\begin{document}
\maketitle

\begin{abstract}
Reliable polyp-size estimation from colonoscopy video requires both useful
image geometry and a defensible link from image evidence to millimetres. We
report preliminary results from a self-supervised monocular depth project that
uses C3VDv2 phantom data as a source-domain guardrail and REALColon clinical
video as an adaptation target. The preferred source baseline is Run D, which
combines SpecularNet and PPS and obtains an AbsRel of 0.1360 and an RMSE of
3.79~mm on the full C3VDv2 HF test. In contrast, unconstrained target-only
adaptation produced forgetting or absolute-scale drift. Freezing depth modules
preserved source outputs, but is a no-forgetting control rather than evidence
of clinical depth adaptation. Teacher-anchored source replay has a promising
sparse clinical-depth signal, but it has not completed target-domain,
qualitative, or scale-aware selection. These findings motivate an
evaluation-first, anchor-gated direct sizing endpoint; they do not demonstrate
metric polyp sizing or clinical 10-mm threshold accuracy.
\end{abstract}

\section{Introduction}

Colonoscopy is a flexible endoscopic procedure used to inspect the large
bowel. Polyps are mucosal lesions that may be removed or followed, and their
estimated diameter can influence documentation, resection planning, and
surveillance pathways. A useful assistive system must operate under wide-angle
optics, changing working distance, specular highlights, blur, debris, bubbles,
insufflation, tissue deformation, and variation in endoscope hardware.

The central technical constraint is metric scale. Under perspective projection,
an image coordinate is proportional to $fX/Z$; scaling scene geometry and
camera translation together leaves the image unchanged. Monocular depth or
multi-view reconstruction can therefore yield useful relative shape without
showing that a measurement is correct in millimetres. Clinical systems that
make metric claims establish a scale channel explicitly, for example through
lens calibration \cite{wang2024endoangel} or an in-scene waterjet reference
\cite{sudarevic2023poseidon}. This short paper reports the project state with
that distinction preserved.

\section{Methods}

\subsection{Research setting}

The implementation extends AF-SfMLearner, a self-supervised monocular
endoscopy depth-and-motion framework that uses appearance flow to mitigate
photometric mismatch \cite{shao2022afsfm}. C3VDv2 phantom data provide dense
depth and pose guardrails. REALColon video provides clinical appearance
diversity but not dense metric depth ground truth. Thus, C3VDv2 evaluates
source-domain geometry and REALColon adaptation must be screened for both
source regression and clinically meaningful scale evidence.

We compare source and adaptation branches documented in the canonical project
record \cite{projectrecord2026}. For every depth-adapting REALColon candidate,
the required evaluation design includes C3VDv2 depth without median scaling for
absolute scale, median-scaled C3VDv2 depth for relative shape, C3VDv2 pose, a
fixed REALColon qualitative panel, and any available sparse clinical gate.
These are selection requirements, not interchangeable endpoints.

\subsection{Interpretation protocol}

We use the following rules to avoid overinterpretation. A frozen module can
prove that its output was preserved, but cannot establish adaptation of that
module. Median-scaled depth measures relative structure and cannot be reported
as metric accuracy. Sparse depth points are not lesion longest-axis labels. A
metric polyp-size claim requires a leakage-checked clinical benchmark with
known label provenance and a validated scale anchor or scope-specific
calibration.

\section{Results}

Table~\ref{tab:findings} summarizes the claim-safe findings. Run D
(\texttt{stage2\_specular\_pps/weights\_best}) is the preferred C3VDv2
phantom-depth baseline. Its full HF-test metrics are AbsRel 0.1360, RMSE
3.79~mm, and $\delta<1.25$ of 0.8318. In the tested ablation, PPS alone was
weaker than the baseline, whereas the SpecularNet--PPS combination was
stronger; this supports a conditional, rather than universal, benefit of PPS.

\begin{table*}[t]
\centering
\footnotesize
\setlength{\tabcolsep}{4pt}
\begin{tabularx}{\textwidth}{@{}p{0.24\textwidth}p{0.30\textwidth}X@{}}
\toprule
\textbf{Finding} & \textbf{Evidence} & \textbf{Safe interpretation} \\
\midrule
Run D is the preferred source baseline. & C3VDv2 HF test: AbsRel 0.1360, RMSE 3.79~mm, $\delta<1.25=0.8318$. & SpecularNet plus PPS is the strongest tested \emph{phantom-depth} configuration; it is not a real-tissue or polyp-size result. \\

PPS depends on specular handling. & PPS-only Run C is weaker than the baseline; the SpecularNet plus PPS Run D is stronger than Run B and Run A. & The tested evidence supports PPS when paired with SpecularNet, not a blanket claim for PPS. \\

Target-only adaptation is unsafe. & Run B and Run D target-only branches failed a planned relative-depth gate; another short full branch retained relative structure but drifted in absolute scale. & Unconstrained adaptation can damage the source prior or metric calibration. \\

Frozen depth is a control. & In \texttt{pose\_transform} and Run B \texttt{teacher\_pose\_transform}, encoder/depth/specular hashes and depth outputs were unchanged. & The controls establish no forgetting for frozen modules, not clinical-domain depth adaptation. \\

Source replay is unselected. & A 2k replay checkpoint improved a sparse PolypSense3D raw gate (MAE 8.5612~mm) while source guardrails remained near baseline. The 3k final checkpoint is clinically interesting but awaits fixed-panel and scale-aware selection. & Replay can protect source performance during trainable-depth adaptation. It does not establish broad clinical metric depth or polyp sizing. \\
\bottomrule
\end{tabularx}
\caption{Claim-safe preliminary results. All adaptation candidates remain
subject to source-depth, source-pose, fixed-panel, and scale-aware gates.}
\label{tab:findings}
\end{table*}

The evidence also identifies what is missing. The sparse clinical-depth signal
measures a limited set of labelled points; it neither measures a lesion's
longest axis nor validates 5-mm or 10-mm decision thresholds. The private
candidate size archive cannot serve as headline evidence until duplicate and
conflicting labels, mixed coordinate conventions, and possible split leakage
are resolved.

\section{Discussion}

The results separate two useful but different outcomes: preserving a source
geometry prior and establishing a clinical measurement endpoint. Run D should
remain the phantom baseline because it is the strongest tested source-domain
configuration. The target-only failures mean that REALColon adaptation should
not be selected by training loss or a single target signal. Conversely, frozen
depth controls are valuable because they localize forgetting, but should not be
misdescribed as successful clinical adaptation.

The direct-sizing literature supports a complementary path. A single-scope,
snare-anchored study has evaluated direct size regression with segmentation and
depth cues \cite{yii2025referencefree}, and MPSE reports direct per-frame
regression with polyp-level aggregation \cite{lawrence2026mpse}. These studies
do not remove scale transfer as a problem, but they show that a useful endpoint
need not be ``reconstruct a point cloud, then measure it.'' The recommended
next track is a short-video direct head with lesion-level aggregation, explicit
quality control, and a metric output only when a validated anchor is visible;
otherwise it should return ordinal probabilities or abstain.

\section{Limitations}

This is an internal preliminary study, not an independent clinical validation.
C3VDv2 is a controlled phantom benchmark, and REALColon lacks dense clinical
depth ground truth. Reported project numbers are traceable to the canonical
experiment index rather than a public checkpoint release. Sparse clinical
points, private candidate labels, and fixed qualitative panels remain
insufficient for metric lesion-size claims. No conclusion here establishes
clinical generalization, regulatory readiness, or real-domain 10-mm threshold
accuracy.

\section{Conclusion}

Run D is the preferred C3VDv2 phantom-depth baseline. Target-only REALColon
adaptation has shown forgetting or scale drift; frozen-depth branches are
no-forgetting controls; and source replay is a guarded but not yet selected
clinical signal. The immediate priority is not a longer adaptation campaign.
It is a target-domain selector and leakage-checked, anchor-annotated clinical
evaluation that can support a direct, clip-level sizing endpoint.

\begingroup
\footnotesize
\bibliographystyle{plain}
\bibliography{references}
\endgroup
\end{document}
